\SourceLeaf{116}
Voraussetzungen, indem man den wahren Werth von \(h\) entweder \(=H\) oder \(=H'\) setzt, so verhalten sich die Wahrscheinlichkeiten, mit welchen sich in denselben die Fehler \(\alpha,\,\beta,\,\gamma,\,\delta\) u.\,s.\,w. erwarten liessen, resp. wie
\[
\begin{gathered}
He^{-HH\alpha\alpha}.He^{-HH\beta\beta}.He^{-HH\gamma\gamma}.\text{ u.\,s.\,w.}\\
\text{zu}\quad H'e^{-H'H'\alpha\alpha}.H'e^{-H'H'\beta\beta}.H'e^{-H'H'\gamma\gamma}.\text{ u.\,s.\,w.}
\end{gathered}
\]
d.\,i. wie
\[
H^m e^{-HH(\alpha\alpha+\beta\beta+\gamma\gamma+\text{ u.\,s.\,w.})}\quad\text{zu}\quad
H'{}^m e^{-H'H'(\alpha\alpha+\beta\beta+\gamma\gamma+\text{ u.\,s.\,w.})}
\]
In demselben Verhältnisse stehen folglich die Wahrscheinlichkeiten, dass \(H\) oder \(H'\) der wahre Werth von \(h\) war, \textit{nach} dem Erfolge jener Fehler (\textit{T. M. C. C.} Art. 176): oder die Wahrscheinlichkeit jedes Werthes von \(h\) ist der Grösse
\[
h^m e^{-hh(\alpha\alpha+\beta\beta+\gamma\gamma+\text{ u.\,s.\,w.})}
\]
proportional. Der \textit{wahrscheinlichste} Werth von \(h\) ist folglich derjenige, für welchen diese Grösse ein Maximum wird, welchen man nach bekannten Regeln
\[
=\surd\frac{m}{2(\alpha\alpha+\beta\beta+\gamma\gamma+\text{ u.\,s.\,w.})}
\]
findet. Der wahrscheinlichste Werth von \(r\) wird folglich
\[
\begin{aligned}
&=\rho\surd\frac{2(\alpha\alpha+\beta\beta+\gamma\gamma+\text{ u.\,s.\,w.})}{m}\\[4pt]
&=0{,}6744897\surd\frac{\alpha\alpha+\beta\beta+\gamma\gamma+\text{ u.\,s.\,w.}}{m}
\end{aligned}
\]
Dies Resultat ist allgemein, \(m\) mag gross oder klein sein.

\Article{4}
Man begreift leicht, dass man von dieser Bestimmung von \(h\) und \(r\) desto weniger berechtigt ist, viele Genauigkeit zu erwarten, je kleiner \(m\) ist. Entwickeln wir daher den Grad von Genauigkeit, welchen man dieser Bestimmung beizulegen hat, für den Fall, wo \(m\) eine grosse Zahl ist. Wir bezeichnen den gefundenen wahrscheinlichen Werth von \(h\)
\[
\surd\frac{m}{2(\alpha\alpha+\beta\beta+\gamma\gamma+\text{ u.\,s.\,w.})}
\]
Kürze halber mit \(H\), und bemerken, dass die Wahrscheinlichkeit, \(H\) sei der
\SourceLeaf{117}
wahre Werth von \(h\), zu der Wahrscheinlichkeit, dass der wahre Werth \(=H+\lambda\) sei, sich verhält, wie
\[
H^m e^{-\frac m2}:(H+\lambda)^m e^{-\frac{m(H+\lambda)^2}{2HH}}
\]
oder wie
\[
1:e^{-\frac{\lambda\lambda m}{HH}(1-\frac13.\frac\lambda H+\frac{1\,\lambda\lambda}{4HH}-\frac15.\frac{\lambda^3}{H^3}\, .\,.\,.\,.)}
\]

Das zweite Glied wird gegen das erste nur dann noch merklich sein, wenn \(\frac\lambda H\) ein kleiner Bruch ist, daher wir uns erlauben dürfen, anstatt des angegebenen Verhältnisses dieses zu gebrauchen
\[
1:e^{-\frac{\lambda\lambda m}{HH}}
\]
Dies heisst nun eigentlich so viel: die Wahrscheinlichkeit, dass der wahre Werth von \(h\) zwischen \(H+\lambda\) und \(H+\lambda+\mathrm{d}\lambda\) liege, ist sehr nahe
\[
=Ke^{-\frac{\lambda\lambda m}{HH}}\mathrm{d}\lambda
\]
wo \(K\) eine Constante ist, die so bestimmt werden muss, dass das Integral
\[
\int Ke^{-\frac{\lambda\lambda m}{HH}}\mathrm{d}\lambda
\]
zwischen den zulässigen Grenzen von \(\lambda\) genommen, \(=1\) werde. Statt solcher Grenzen ist es hier, wo wegen der Grösse von \(m\) offenbar
\[
e^{-\frac{\lambda\lambda m}{HH}}
\]
unmerklich wird, sobald \(\frac\lambda H\) aufhört ein kleiner Bruch zu sein, erlaubt, die Grenzen \(-\infty\) und \(+\infty\) zu nehmen, wodurch
\[
K=\frac1H\surd\frac m\pi
\]
wird. Mithin ist die Wahrscheinlichkeit, dass der wahre Werth von \(h\) zwischen \(H-\lambda\) und \(H+\lambda\) liege,
\[
=\theta\bigl(\tfrac\lambda H\surd m\bigr)
\]
also jene Wahrscheinlichkeit \(=\frac12\), wenn
\[
\frac\lambda H\surd m=\rho\quad\text{ist.}
\]
